% ============================================================================
\chapter{Mjerenje zamjene konteksta i pore\dj{}enje sa simulatorom}
\label{ch:mjerenje}
% ============================================================================

U ovom poglavlju opisuje se mjerenje vremena zamjene konteksta operativnog
sistema na stvarnom hardveru, te njegovo pore\dj{}enje sa ranije dobijenim
vrijednostima iz simulatora. Najprije se obja\v{s}njava metodologija mjerenja i
razlog za odre\dj{}ene odluke pri njenoj izradi. Zatim se prikazuju izmjereni
rezultati na hardveru. Na kraju se ti rezultati pa\v{z}ljivo pore\dj{}uju sa
rezultatima iz simulatora, uz jasno isticanje ograni\v{c}enja takvog pore\dj{}enja.
Cilj poglavlja je dati prvu mjerljivu karakterizaciju operativnog sistema na
ovoj platformi, bez preuveli\v{c}avanja zaklju\v{c}aka.

\section{Metodologija mjerenja}
\label{sec:mjerenje-metoda}

Mjerenje oponasa test zamjene konteksta iz simulatorskog dijela istra\v{z}ivanja
\citep{thesisqemu}. Jedna iteracija obuhvata dvije zamjene konteksta, odlazak iz
prvog u drugi zadatak i povratak nazad. Vrijeme se mjeri slobodnim hardverskim
broja\v{c}em koji radi na frekvenciji od dvadeset sedam megaherca. Za razliku od
simulatorskog testa, koji koristi dva zadatka istog prioriteta i jedan semafor,
\v{c}ije pona\v{s}anje zavisi od dijeljenja vremena, ovdje su kori\v{s}tena dva
semafora i zadatak odgovara\v{c} vi\v{s}eg prioriteta. Time svaka predaja semafora
sigurno izaziva mjerenu zamjenu konteksta. Mjerena veli\v{c}ina je ista, broj
takta za dvostruku zamjenu konteksta, a jednostruka zamjena dobija se dijeljenjem
sa dva.

\section{Rezultati na hardveru}
\label{sec:mjerenje-rezultati}

Mjereno je pet stotina iteracija, a rezultat je bio ponovljiv kroz vi\v{s}e
upisa u plo\v{c}u. Dobijene vrijednosti za dvostruku zamjenu konteksta iznose
sedam hiljada \v{s}ezdeset jedan takt kao najmanja, sedam hiljada sedamdeset
devet kao prosje\v{c}na i sedam hiljada devet stotina osamdeset dva kao najve\'ca
vrijednost. To odgovara prosje\v{c}nom vremenu od oko dvije stotine \v{s}ezdeset
dvije mikrosekunde. Jednostruka zamjena konteksta iznosi oko tri hiljade pet
stotina trideset takta, odnosno oko sto trideset jednu mikrosekundu. Mali raspon
izme\dj{}u najmanje i najve\'ce vrijednosti pokazuje da nijedan takt
raspore\dj{}iva\v{c}a nije upao u prozor mjerenja, pa su vrijednosti \v{c}iste.

Visok broj takta po zamjeni konteksta posljedica je toga \v{s}to je PicoRV32
vi\v{s}eciklusna jezgra na relativno niskoj frekvenciji, kako je opisano u
poglavlju \ref{ch:platforma}. Zamjena konteksta pohranjuje i obnavlja cijeli
skup cjelobrojnih registara, uz dodatni tro\v{s}ak rada sa semaforima.

\section{Pore\dj{}enje sa simulatorom}
\label{sec:mjerenje-poredjenje}

Rezultati iz simulatora, dobijeni emulacijom RISC-V sistema, daju za dvostruku
zamjenu konteksta najmanju vrijednost od devet hiljada sto osamdeset osam takta,
prosje\v{c}nu od dvadeset tri hiljade sedam stotina \v{c}etrdeset sedam i
najve\'cu od gotovo \v{c}etiri miliona takta. Pore\dj{}enje je dato u tabeli
\ref{table:poredjenje}.

\begin{table}[H]
\caption{Pore\dj{}enje dvostruke zamjene konteksta na hardveru i u simulatoru (u taktovima).}
\label{table:poredjenje}
\begin{center}
\begin{tabular}{lcc}
\hline
Veli\v{c}ina & Hardver (27 MHz) & Simulator (25 MHz) \\ \hline
Najmanja & 7061 & 9188 \\
Prosje\v{c}na & 7079 & 23747 \\
Najve\'ca & 7982 & 3895040 \\ \hline
\end{tabular}
\end{center}
\end{table}

Pore\dj{}enje treba tuma\v{c}iti oprezno. Hardverska vrijednost dobijena je
stvarnim broja\v{c}em takta, dok simulatorska vrijednost dolazi iz emuliranog
broja\v{c}a koji nije takovno ta\v{c}an u odnosu na stvarni \v{c}ip. Zbog toga se
ne pore\dj{}uju dva ista takta, nego model iz simulatora i stvarni hardver.
Najpo\v{s}tenija ta\v{c}ka pore\dj{}enja je najmanja vrijednost, jer nije optere\'cena
smetnjama, i tu su vrijednosti istog reda veli\v{c}ine. Prosje\v{c}na i najve\'ca
vrijednost u simulatoru znatno su ve\'ce zbog ka\v{s}njenja koja unosi
periodi\v{c}no \v{c}ekanje u testu i takt raspore\dj{}iva\v{c}a. Zbog svega
navedenog, ne iznosi se tvrdnja o ta\v{c}nom odnosu brzina, nego se zaklju\v{c}uje
da je izmjerena hardverska zamjena konteksta po redu veli\v{c}ine u skladu sa
\v{c}istom procjenom iz simulatora.

\section{Zaklju\v{c}ak poglavlja}
\label{sec:mjerenje-zakljucak}

U poglavlju je dato prvo mjerenje vremena zamjene konteksta operativnog sistema
na stvarnom hardveru. Opisana je metodologija, prikazani su ponovljivi rezultati
i izvr\v{s}eno je oprezno pore\dj{}enje sa simulatorom. Pokazano je da su \v{c}iste
vrijednosti istog reda veli\v{c}ine, ali i da razlike u izvoru takta i uslovima
mjerenja ne dopu\v{s}taju izvo\dj{}enje ta\v{c}nog odnosa brzina. Time je
zaokru\v{z}ena mjerljiva karakterizacija sistema na ovoj platformi.
